\section{Análise no Domínio da Frequência}

\begin{frame}{Diagrama de Bode}
  \begin{columns}
    \column{0.64\textwidth}
    \scriptsize
    \begin{equation}
      g_{3,3}(s) \approx \frac{- 5{,}6 s^{3} - 108{,}7 s^{2} + 1249{,}8 s + 15969{,}5}{s^{4} + 106{,}2 s^{3} + 3688{,}8 s^{2} + 29427{,}6 s + 59584{,}6}
    \end{equation}
    \begin{equation}
      g_{3,3}(jw) \approx \frac{5{,}6 j w^{3} + 108{,}7 w^{2} + 1249{,}8 j w + 15969{,}5}{w^{4} - 106{,}2 j w^{3} - 3688{,}8 w^{2} + 29427{,}6 j w + 59584{,}6}
    \end{equation}
    \column{0.475\textwidth}
    \begin{itemize}
      \item A função de transferência é avaliada em função da frequência de excitação (jw).
      \item Comportamento do sistema a diferentes frequências de excitação.
      \item Possibilidade de decompor o sistema e analisar cada parcela separadamente.
      \item Auxílio no projeto de Filtros.
    \end{itemize}
  \end{columns}
\end{frame}

\begin{frame}{Diagrama de Bode}
  \begin{columns}
    \column{0.5\textwidth}
    \begin{figure}
      \centering
      \includegraphics[width=\textwidth]{figuras/T3_FR_bode_freq.png}
      \caption{Diagrama de bode da função de transferência que relaciona $F_R/T_3$ com marcações nas frequências de canto e natural.}
    \end{figure}
    \column{0.5\textwidth}
    \begin{block}{Influência dos Polos e Zeros ($n$ = ordem)}
      \begin{itemize}
        \item Polo: $-20n$ dB/década | $-90n^\circ$
        \item Zero: $+20n$ dB/década | $+90n^\circ$
      \end{itemize}
    \end{block}
    \begin{itemize}
      \item Ganho estático do sistema é mantido até aproximadamente 1 rad/h.
      \item Nesse sistema as frequências de canto de polos e zeros estão muito próximas, sobrepondo seus efeitos.
      \item O sistema é de no mínimo 3ª ordem, com um zero no semi-plano direito.
    \end{itemize}
  \end{columns}
\end{frame}

% \begin{frame}{Diagrama de Bode}
%   \begin{columns}
%     \column{0.5\textwidth}
%     \begin{figure}
%       \centering
%       %MUDAR A IMAGEM PARA A QUE ESTÁ COM AS FREQUÊNCIAS DE CANTO
%       \includegraphics[width=\textwidth]{figuras/T3_FR_bode_freq.png}
%       \caption{Diagrama de bode com as frequências de canto dos polos e a frequência natural explícitas.}
%     \end{figure}
%     \column{0.5\textwidth}
%     \begin{table}
%       \caption{Frequência de canto de cada polo real}
%       \begin{tabular}{cc}
%         \hline
%         Símbolo & Frequência de Canto / rad/h \\
%         \hline
%         $w_{1}$ & 3 \\
%         $w_{2}$ & 7 \\
%         \hline
%       \end{tabular}
%     \end{table}
%     \begin{table}
%       \caption{Frequências natural da segunda ordem}
%       \begin{tabular}{cc}
%         \hline
%         Símbolo & Frequência Natural / rad/h \\
%         \hline
%         $w_{3}$ & 50 \\
%         \hline
%       \end{tabular}
%     \end{table}
%   \end{columns}
% \end{frame}

% \begin{frame}{Diagrama de Bode}
%   \begin{columns}
%     \column{0.5\textwidth}
%     Parcelas de Primeira Ordem:
%     \begin{block}{Constante de Tempo}
%       \centering
%       $\tau = \frac{1}{w_{canto}}$
%     \end{block}
%     \begin{table}
%       \caption{Constante de Tempo}
%       \begin{tabular}{cc}
%         \hline
%         $w_{canto}$ / rad/h & $\tau$ / h \\
%         \hline
%         3 & $0{,}33$ \\
%         7 & $0{,}14$ \\
%         \hline
%       \end{tabular}
%     \end{table}
%     \column{0.5\textwidth}
%     Parcela de Segunda Ordem:
%     \begin{itemize}
%       \item O termo de segunda ordem possui polos complexos conjugados.
%       \item Ausência de pico de ressonância no diagrama.
%     \end{itemize}
%     \begin{block}{Fator de Amortecimento}
%       \centering
%       $\frac{1}{\sqrt{2}}<\xi<1$
%     \end{block}
%     \begin{block}{Constante de Tempo}
%       \centering
%       $\tau = \frac{1}{\xi wn}$ \\
%       $0,03$ h $< \tau < 0,02$ h
%     \end{block}
%   \end{columns}
% \end{frame}

\begin{frame}{Projeto de Filtro}
  \begin{columns}
    \column{0.55\textwidth}
    \begin{figure}
      \centering
      \includegraphics[width=\textwidth]{figuras/T3_FR_noise.png}
      \caption{Resposta ao degrau unitário em desvio de $F_R/T_3$ do sistema não linear, com ruído gaussiano.}
    \end{figure}
    \column{0.45\textwidth}
    Ruído gaussiano com média zero e desvio padrão $\sigma = 0{,}02.$
    \begin{block}{Objetivo}
      Reduzir o efeito do ruído preservando a dinâmica do sistema.
    \end{block}
  \end{columns}
\end{frame}

\begin{frame}{Projeto de Filtro}
  \begin{columns}
    \column{0.5\textwidth}
    \begin{figure}
      \centering
      \includegraphics[width=\textwidth]{figuras/T3_FR_fft.png}
      \caption{Espectro de frequência da resposta ruidosa do sistema.}
    \end{figure}
    \column{0.5\textwidth}
    Para reduzir o risco de atenuação de componentes relevantes do sinal, adotou-se
    $$\omega_{\mathrm{corte}} = 40\ \mathrm{rad/h}$$
    \vspace{-2mm}
    \begin{block}{Filtro de 1ª Ordem}
      $$F_1(s)=\frac{1}{\frac{1}{40}s+1}$$
    \end{block}
    \begin{block}{Filtro de 2ª Ordem}
      $$F_2(s)=\frac{40^2}{s^2+\frac{2\cdot40}{\sqrt{2}}s+40^2}$$
    \end{block}
  \end{columns}
\end{frame}

\begin{frame}{Projeto de Filtro}
  \begin{columns}
    \column{0.5\textwidth}
    \begin{figure}
      \centering
      \includegraphics[width=\textwidth]{figuras/T3_FR_bode_filter.png}
      \caption{Comparação do diagrama de bode dos dois filtros.}
    \end{figure}
    \column{0.5\textwidth}
    \begin{itemize}
      \item O filtro de 2ª ordem oferece o dobro de atenuação na banda de rejeição do que o de 1ª ordem.
      \item Entretanto, oferece uma maior defasagem, gerando atraso em parte da resposta do sinal.
    \end{itemize}
  \end{columns}
\end{frame}

\begin{frame}{Comparação dos Filtros}
  \begin{columns}
    \column{0.7\textwidth}
    \begin{figure}
      \centering
      \includegraphics[width=\textwidth]{figuras/T3_FR_compare.png}
      \caption{Comparação da eficácia dos dois filtros.}
    \end{figure}
    \column{0.3\textwidth}
    \begin{itemize}
      \item Ambos filtram o sinal de forma eficaz.
      \item O filtro de segunda ordem apresenta uma atenuação ligeiramente superior, à custa de um aumento no defasamento da resposta.
    \end{itemize}
  \end{columns}
\end{frame}
